\chapter{Testing the tests}\label{ch:robust}

The third notebook, \texttt{climate\_data\_robustness.ipynb}, removes the shortcuts of the second. It reads the
daily archive the original study used and compares it with the monthly one, station by station; it builds a
homogenisation method from scratch (the Standard Normal Homogeneity Test, tested first on made-up data) and compares
the result with NOAA's and with Berkeley Earth's; it runs the energy-balance model with 135 combinations of
sensitivity and timescales; it reads ten runs each of two climate models to measure chance; and it computes the
trend for every start year and window length. Run the second notebook first, or let this one download what it
needs.

\section{The plan}
\input{generated/nb_climate_data_robustness_1}
\section{Daily data}
\input{generated/nb_climate_data_robustness_2}
\section{Our own homogenisation}
\input{generated/nb_climate_data_robustness_3}
\section{The energy-balance model, 135 ways}
\input{generated/nb_climate_data_robustness_4}
\section{Many runs of one model}
\input{generated/nb_climate_data_robustness_5}
\section{Every window, not just four}
\input{generated/nb_climate_data_robustness_6}
\section{What this notebook adds}
\input{generated/nb_climate_data_robustness_7}
\section{Recording what the run used}
\input{generated/nb_climate_data_robustness_8}

\section{Where to go next}
\begin{itemize}
\item Run the homogenisation on all 1,155 of the study's stations (about ten minutes) and compare US and non-US averages.
\item Add a metadata step: NOAA's station-history database (HOMR) records when US cooperative observers changed their reading time; use it to correct the change from afternoon to morning readings and see how much of NOAA's adjustment it explains.
\item Try a different break test (for example a penalised likelihood or the Pettitt test) and compare the breaks found.
\end{itemize}
